%! TEX root = ../ag-19-20.tex

\chapter{Examen 2018--2019}

\section{Seria A}

\textbf{Numărul 1}

1. În spațiul vectorial $ \RR^3 $ se consideră mulțimea:
\[
  V = \{ (a, b, c) \in \RR^3 \mid a - 2b + 3c = 0 \}.
\]

Să se determine $V^\perp$ relativ la produsul scalar euclidian și o
bază ortonormată în $ \RR^3 $ diferită de baza canonică.

\vspace{1cm}

2. Să se aducă la forma canonică și să se reprezinte grafic conica:
\[
  f(x, y) = 9x^2 + 4xy - 3y^2 + 4x - 14y - 7 = 0.
\]


\vspace{1cm}

3. Să se rezolve ecuațiile diferențiale:
\begin{enumerate}[(a)]
  \item $ xy' - y + 2x^2y^2 = 0 $;
  \item $ y = x(1 + y') + (y')^2 $;
  \item $ (3x^2 + 2) y' + xy^2 = 0 $;
  \item $ (2x - y + 2) dx + (-x + y + 1) dy = 0 $.
\end{enumerate}

\vspace{1cm}

4. Să se rezolve sistemul:
\[
  \begin{cases}
    x' &= -3x - y + e^t \\
    y' &= x - y
  \end{cases}
\]
cu condițiile inițiale $ x(0) = 1, y(0) = -1 $ și să se calculeze
$ e^{At}$, unde $ A $ este matricea sistemului.

\vspace{1cm}

5.(a) Să se determine ecuația conului cu vîrful în $ V(2, -1, 0) $ și
determinat de curba\footnotemark:
\[
  (\gamma) : \begin{cases}
    y^2 + z^2 &= 2 \\
    x &= 1
  \end{cases}
\]

\footnotetext{Pentru ecuația conului și a altor cuadrice nestudiate
    la seminar, consultați notițele Prof.\ Niță sau, de exemplu,
\href{https://www.deliu.ro/pluginfile.php/109/mod_resource/content/1/cap02.pdf}{acest document}.}

(b) Se consideră spațiul euclidian $ \RR_4[X] $ și un endomorfism
$ f : \RR_4[X] \to \RR_4[X] $. Să se arate că:
\[
  \dr{Im}f \simeq (\dr{Ker}f^\ast)^\perp.
\]



\textbf{Numărul 2}

1. În spațiul euclidian $ \RR^3 $ se consideră mulțimea:
\[
  V = \{(a, b, c) \in \RR^3 \mid 2a - b + 3c = 0 \}.
\]

Să se determine $ V^\perp $ și o bază ortonormată în $ \RR^3 $, diferită
de baza canonică.

\vspace{1cm}


2. Să se aducă la forma canonică și să se reprezinte grafic conica:
\[
  f(x, y) = 9x^2 + 4xy + 3y^2 + 16x + 12y - 36 = 0.
\]

\vspace{1cm}

3. Să se rezvole ecuațiile diferențiale:
\begin{enumerate}[(a)]
  \item $ xy' = 4y + x^2 \sqrt{y} $;
  \item $ yy' = 2x(y')^2 + 1 $;
  \item $ (x^2 + 1)y' + 2xy^2 = 0 $;
  \item $ (2x + y + 1)dx + (x + y - 1) dy = 0 $.
\end{enumerate}


\vspace{1cm}

4. Să se rezolve sistemul:
\[
  \begin{cases}
    x' &= 2x + y \\
    y' &= -x + 4y + 27t
  \end{cases}
\]
cu valorile inițiale $ x(0) = -1, y(0) = 1 $ și să se calculeze
$ e^{At} $, unde $ A $ este matricea sistemului.

\vspace{1cm}

5.(a) Să se determine ecuația cilindrului cu generatoarele paralele
cu vectorul $ \vec{v} = \vec{i} - \vec{j} $ și curba directoare:
\[
  (\gamma): %
  \begin{cases}
    x^2 - y^2 + z^2 &= 0 \\
    x &= 1
  \end{cases}
\]

(b) Se consideră spațiul euclidian $ \RR_4[X] $ și $f : \RR_4[X] \to \RR_4[X] $
un endomorfism. Să se arate că:
\[
  \dr{Ker}f \simeq (\dr{Im}f^\ast)^\perp.
\]

%%%%%%%%%%%%%%%%%%%%%%%%%%%%%%%%%%%%%%%%%%%%%%%%%%%%%%%%%%%%%%%%%%%%%%
\section{Seria D}

\textbf{Observație: Numărul 1 = (i), numărul 2 = (ii) etc.}

1. (a) Folosind metoda Gram-Schmidt, ortonormați familiile
de vectori:
\begin{enumerate}[(i)]
  \item $ B = \{ v_1 = (0, 1, 1), v_2 = (1, 1, 3), v_3 = (4, 2, 0) \} $;
  \item $ B = \{ v_1 = (2, 0, 0), v_2 = (-3, 1, 1), v_3 = (2, -2, 5) \} $;
  \item $ B = \{ v_1 = (1, 0, 1), v_2 = (0, 2, 4), v_3 = (1, 1, 1) \}$;
  \item $ B = \{ v_1 = (1, 1, 0), v_2 = (-2, 0 3), v_3 = (3, 1, 1) \}$.
\end{enumerate}

(b) (i) Fie $ U = \dr{Sp}\{u_1 = (1, 3, 0, 1), u_2 = (-1, 0, 1, 1) \} \leq \RR^4 $.

Să se determine complementul ortogonal $ U^\perp $ și o bază a acestuia.

(ii) Fie $ \alpha \in (0, 2\pi) $ fixat. Arătați că endomorfismul:
\[
  T(x_1, x_2) = (x_1 \sin \alpha + x_2 \cos \alpha, x_1 \cos \alpha - x_2 \sin \alpha)
\]
este ortogonal.

(iii) Fie matricea $ A = \begin{pmatrix} 1 & 0 & 3 \\ -1 & 2 & 0 \\ 0 & 0 & 1 \end{pmatrix} $.
Folosind teorema Cayley-Hamilton, aflați $ A^{-1} $.

(iv) Fie vectorii $ u = (1, 2, 1), v = (2, 3, 0), w = (0, 1, 1) $.

Să se calculeze aria paralelogramului determinat de $ u $ și $ v $ și volumul
paralelipipedului determinat de $ u, v $ și $ w $.

\vspace{1cm}

2. Să se identifice, aducă la forma canonică și reprezinte grafic conica:
\begin{enumerate}[(i)]
  \item $ x^2 - 2xy + y^2 - 6x - 2y + 9 = 0 $;
  \item $ x^2 - xy + y^2 - 5x + y - 2 = 0 $;
  \item $ 5x^2 + 6xy + 5y^2 - 16x - 16y - 16 = 0 $;
  \item $ 3x^2 - 10xy + 3y^2 + 4x + 4y + 4 = 0 $.
\end{enumerate}

\vspace{1cm}

3. Să se identifice și să se rezolve următoarele ecuații diferențiale
de ordinul întîi:
(i) \begin{enumerate}[(a)]
  \item $ y' = \dfrac{x^2}{y} $;
  \item $ y' = \dfrac{y}{x} + x $;
  \item $ y = xy' + \sin y' $;
  \item $ 3x^2y dx - (\ln y + x^3) dy = 0 $.
\end{enumerate}

(ii) \begin{enumerate}[(a)]
  \item $ y' = x^2 y $;
  \item $ y' = \dfrac{x^2 + y^2}{2xy} $;
  \item $ y' = - \dfrac{1}{2x^2} y^2 + \dfrac{2}{x} y $;
  \item $ y = -xy' + \ln y' $.
\end{enumerate}

(iii)
\begin{enumerate}[(a)]
  \item $ y' = 2xy^2 $;
  \item $ y' = \dfrac{x - y}{x + y}, x + y \neq 0 $;
  \item $ y = xy' - e^{y'} $;
  \item $ \dfrac{y}{x} dx + (y^3 - \ln x) dy = 0 $.
\end{enumerate}

(iv)
\begin{enumerate}[(a)]
  \item $ y' = \dfrac{4x}{y^2} $;
  \item $ y' = \dfrac{2xy + x}{1 + x^2} $;
  \item $ (x + y + 1) dx + (x - y^3 + 3)dy = 0 $;
  \item $ y = \dfrac{{y'}^2 + 4}{2y'} $.
\end{enumerate}


\vspace{1cm}

4. Folosind matricea $ e^{At} $, să se rezolve următoarele sisteme:
\begin{enumerate}[(i)]
  \item $ \begin{cases}
      x' &= y + z \\
      y' &= 2y \\
      z' &= -2x + y + 3z
    \end{cases} $;
  \item $ \begin{cases}
      x' &= 2x \\
      y' &= 4x + y \\
      z' &= -2x + z
    \end{cases};
    $
  \item $ \begin{cases}
      x' &= -2z \\
      y' &= x + 2y + z \\
      z' &= x + 3z
    \end{cases} $;
  \item $ \begin{cases}
      x' &= x - 2z \\
      y' &= y + 4x \\
      z' &= 2z
    \end{cases} $.
\end{enumerate}


%%%%%%%%%%%%%%%%%%%%%%%%%%%%%%%%%%%%%%%%%%%%%%%%%%%%%%%%%%%%%%%%%%%%%%
\section{Seria E}


\textbf{Observație: Nr.\ 1 = (a), Nr.\ 2 = (b).}

1. Determinați forma canonică Jordan pentru matricea:
\begin{enumerate}[(a)]
  \item $ A = \begin{pmatrix}
      3 & 2 & 2 \\
      2 & 2 & 1 \\
      -6 & -5 & -4
    \end{pmatrix}
    $
  \item $ A = \begin{pmatrix}
      1 & -1 & -1 \\
      -3 & -4 & -3 \\
      4 & 7 & 6
    \end{pmatrix}
    $
\end{enumerate}

\vspace{1cm}

2. Aduceți la forma canonică următoarea conică:
\begin{enumerate}[(a)]
  \item $ \Gamma: x_1^2 - 12x_1x_2 - 4x_2^2 + 12x_1 + 8x_2 + 5 = 0 $;
  \item $ \Gamma: 3x_1^2 - 10x_1x_2 + 3x_2^2 + 4x_1 + 4x_2 + 4 = 0 $.
\end{enumerate}

\vspace{1cm}

3. Rezolvați ecuația diferențială liniară:
\begin{enumerate}[(a)]
  \item $ y' = \dfrac{1}{x} \cdot y - \dfrac{\ln x}{x}, x > 0 $;
  \item $ y' = -\dfrac{y}{x} + \dfrac{e^x}{x}, x > 0 $.
\end{enumerate}

\vspace{1cm}

4. Rezolvați ecuația diferențială, cu factorul integrant indicat:
\begin{enumerate}[(a)]
  \item $ (x + y^2)dx + y(1 - x)dy = 0, \mu = \mu(x) $;
  \item $ 3x^2 y dx - (x^3 + \ln y) dy = 0, \mu = \mu(y) $.
\end{enumerate}

\vspace{1cm}

5. Aflați forma generală a soluției sistemului de ecuații diferențiale:
\begin{enumerate}[(a)]
  \item $ \begin{cases}
      y_1' &= y_1 - 5y_2 \\
      y_2' &= 2y_1 - y_2
    \end{cases}
    $;
  \item $ \begin{cases}
      y_1' &= 4y_1 - 3y_2 \\
      y_2' &= 3y_1 + 4y_2
    \end{cases}
    $.
\end{enumerate}



\textbf{Observație: Nr.\ 1 = (a), Nr.\ 2 = (b).}

1. Determinați forma canonică Jordan a matricei:
\begin{enumerate}[(a)]
  \item $ A = %
    \begin{pmatrix}
      2 & -3 & 4 \\
      4 & -6 & 8 \\
      6 & -7 & 8
    \end{pmatrix} $
  \item $ A = %
    \begin{pmatrix}
      0 & -1 & 0 \\
      -1 & -2 & 2 \\
      0 & -1 & 0
    \end{pmatrix} $.
\end{enumerate}

\vspace{1cm}

2. Aduceți la forma canonică:
\begin{enumerate}[(a)]
  \item $ \Gamma: 3x_1^2 + 4x_1x_2 - 2x_1 + 4x_2 - 3 = 0 $;
  \item $ \Gamma: 5x^2 + 8xy + 5y^2 - 18x - 18y + 9 = 0 $.
\end{enumerate}

\vspace{1cm}

3. rezolvați ecuația diferențială liniară:
\begin{enumerate}[(a)]
  \item $ y' = -2y + x^2 + 2x $;
  \item $ y' = \dfrac{2y}{x} + \dfrac{3}{x^2} $.
\end{enumerate}

\vspace{1cm}

4. Rezolvați ecuația diferențială de tip Bernoulli:
\begin{enumerate}[(a)]
  \item $ y' = 2xy + x^3 y^{\frac{1}{2}} $;
  \item $ y' = -\dfrac{1}{x} y + \dfrac{\ln x}{x} y^2 $.
\end{enumerate}

\vspace{1cm}

5. Să se rezolve sistemele de ecuații diferențiale:
\begin{enumerate}[(a)]
  \item $ \begin{cases}
      y_1' &= 2y_1 + y_2 \\
      y_2' &= -y_1 + 4y_2
    \end{cases} $
  \item $ \begin{cases}
      y_1' &= 5y_1 - y^2 \\
      y_2' &= y_1 + 3y_2
    \end{cases} $.
\end{enumerate}


